% Vietnamese translation for OI Public Library
% Source-PDF: IOI中国国家候选队论文/2004/朱泽园.pdf
\documentclass[11pt]{article}
\usepackage{oipl-vn}

\title{Thuật toán khớp nhiều xâu và những gợi ý từ nó}
\author{Zhu Zeyuan\\Trường Ngoại ngữ Nam Kinh}
\date{}

\begin{document}
\maketitle

\origtitle{Multiple string matching algorithms and their implications}
\sourcepath{IOI China candidate team papers / 2004 / Zhu Zeyuan.pdf}

\section*{Từ khóa}

Xâu mẫu; cây tiền tố từ; cây hậu tố; khớp xâu.

\section*{Trích yếu}

Xử lý xâu có vị trí quan trọng trong ứng dụng thực tế. Nó trông có vẻ đơn
giản, nhưng khi nghiên cứu sâu hơn, nhiều tư tưởng tinh tế xuất hiện trong đó
và độ khó cũng trở nên rất lớn. Vì vậy, các kỳ thi tin học thường dùng bài
toán xử lý xâu để rèn luyện năng lực sáng tạo của thí sinh.

Bài viết nêu bài toán và phân tích trực tiếp ban đầu ở Chương 1. Chương 2,
3 và 5 lần lượt giới thiệu ba công cụ chuẩn bị: thuật toán khớp mẫu KMP, cấu
trúc cây tiền tố từ do tác giả xây dựng, và thuật toán cây hậu tố. Dựa trên tư
tưởng cốt lõi của KMP, Chương 4 đưa ra một thuật toán tuyến tính cho bài toán
khớp nhiều xâu. Bài viết không dừng ở độ phức tạp tuyến tính; Chương 6 tiếp
tục đề xuất một thuật toán có hiệu năng trung bình tốt hơn. Cuối cùng, Chương
7 phân tích cách hình thành thuật toán và nâng tư tưởng ấy lên mức phương
pháp luận.

\tableofcontents

\section{Đặt vấn đề}

\subsection{Mô tả bài toán}

Khớp nhiều xâu nghĩa là: cho một số xâu mẫu, trong một văn bản chỉ gồm 26 chữ
cái thường \texttt{a} đến \texttt{z}, tìm vị trí xuất hiện đầu tiên của bất kỳ
xâu mẫu nào. Cụ thể, cho $m$ xâu mẫu
\[
  P_1[1..L_1], P_2[1..L_2], \ldots, P_m[1..L_m]
\]
và văn bản $T[1..n]$. Với ràng buộc
\[
  \sum L \le 100K,\qquad m\le 1000,\qquad n\le 900K,
\]
ta cần tìm số nguyên nhỏ nhất $s\in[1,n]$ sao cho tồn tại $a\in[1,m]$ thỏa
\[
  \forall x\in[1,L_a],\qquad T[s+x-1]=P_a[x].
\]

Nếu các mẫu là \texttt{cdefg} và \texttt{efg}, còn văn bản là
\texttt{abcdefgh}, yêu cầu ``tìm tất cả vị trí xuất hiện của tất cả mẫu'' dễ
làm mục tiêu trở nên không rõ ràng khi chỉ cần một câu trả lời. Vì vậy ta
thường chuyển sang tìm ``vị trí đầu tiên''. Các thuật toán bên dưới vẫn có thể
liệt kê mọi lần xuất hiện mà không đổi bậc độ phức tạp.

Một ứng dụng thực tế là công cụ tìm kiếm có từ khóa logic. Trong sinh học, khi
tìm một vài mẫu có thể là biến dị trong chuỗi DNA, ta cũng gặp bài toán cùng
loại. Do đó, giải bài toán này bằng thuật toán hiệu quả có thể cải thiện đáng
kể hiệu suất trong nhiều lĩnh vực.

\subsection{Những ý tưởng ban đầu}

Cách đơn giản nhất là xét từng vị trí rồi thử từng mẫu:

\begin{verbatim}
for i <- 1 to n do
  for j <- 1 to m do
    if i + L[j] - 1 <= n then
      if T[i..i+L[j]-1] = P[j][1..L[j]] then
        output i and stop
\end{verbatim}

Thuật toán này kiểm tra các vị trí theo thứ tự tăng dần. Trong trường hợp xấu
nhất, độ phức tạp thời gian là
\[
  O\left(n\sum L\right).
\]

Một cải tiến tự nhiên hơn là tìm lần xuất hiện đầu tiên của từng mẫu trong
văn bản:

\begin{verbatim}
for i <- 1 to m do
  X[i] <- first occurrence of P[i] in T, or infinity if absent
output min(X)
\end{verbatim}

Vị trí xuất hiện đầu tiên của $P_i$ trong $T$ có thể tìm bằng KMP trong
$O(n+L_i)$, do đó tổng độ phức tạp là
\[
  O\left(nm+\sum L\right).
\]

Hai phương pháp này vẫn quá yếu trước quy mô dữ liệu đã nêu. Ta cần cố gắng
tìm một thuật toán tuyến tính, hoặc còn tốt hơn. Trước hết ta giới thiệu KMP,
thuật toán khớp một xâu, và một cấu trúc dữ liệu chuẩn bị là cây tiền tố từ.

\section{Thuật toán Knuth--Morris--Pratt}

Thuật toán này được D. E. Knuth, J. H. Morris và V. R. Pratt độc lập phát
hiện, thường gọi tắt là KMP.

\subsection{Định nghĩa}

Cho mẫu $P[1..m]$ và văn bản $T[1..n]$. Ta cần tìm mọi số nguyên
$s\in[1,n-m+1]$ sao cho
\[
  \forall x\in[1,m],\qquad T[s+x-1]=P[x].
\]

\subsection{Hàm tiền tố của mẫu}

Với $1\le i\le m$, định nghĩa hàm tiền tố
\[
  \pi(i)=\max\{j\mid 0\le j<i,\; P[1..j]=P[i-j+1..i]\}.
\]
Nói cách khác, $\pi(i)$ là độ dài biên dài nhất của $P[1..i]$, tức một tiền tố
đồng thời là hậu tố thật sự của đoạn này.

Hàm này có thể được tiền xử lý bằng đoạn sau:

\begin{verbatim}
k <- 0
pi[1] <- 0
for i <- 2 to m do
  while k > 0 and P[k+1] != P[i] do
    k <- pi[k]
  if P[k+1] = P[i] then
    k <- k + 1
  pi[i] <- k
\end{verbatim}

Vì $k$ tăng nhiều nhất $m-1$ lần, số lần thực hiện phép gán
\texttt{k <- pi[k]} cũng không vượt quá $m-1$. Vì vậy đoạn tiền xử lý có độ
phức tạp $O(m)$.

\subsection{Thuật toán chính của KMP}

Hàm tiền tố ghi lại thông tin tự khớp sau khi dịch của chính mẫu. Khi gặp ký
tự không khớp, ta có thể dựa vào hàm tiền tố để điều chỉnh vị trí mẫu thay vì
so sánh lại từ đầu.

Ví dụ, nếu đang khớp được một đoạn dài $5$ và $\pi(5)=3$, mẫu có thể dịch
sang phải $5-3=2$ ô rồi tiếp tục so sánh. Nếu vẫn không khớp, ta lại tiếp tục
dùng $\pi$ để dịch.

Tương tự cách tính mảng tiền tố, thuật toán khớp chính có dạng:

\begin{verbatim}
q <- 0
for i <- 1 to n do
  while q > 0 and P[q+1] != T[i] do
    q <- pi[q]
  if P[q+1] = T[i] then
    q <- q + 1
  if q = m then
    output an occurrence ending at i
    q <- pi[q]
\end{verbatim}

Giống phần tính hàm tiền tố, số lần thực hiện \texttt{q <- pi[q]} không vượt
quá $n$. Vì vậy toàn bộ KMP chạy trong
\[
  O(n+m).
\]

\section{Cây tiền tố từ}

\subsection{Định nghĩa trie}

Cây tiền tố từ thực chất là một trie. Trong trường hợp lý tưởng, đó là một cây
26 phân vô hạn; 26 cạnh đi xuống từ mỗi nút được gắn nhãn lần lượt là
\texttt{a}, \texttt{b}, \ldots, \texttt{z}.

Với mỗi nút $p$, ghép các chữ cái trên đường đi từ gốc đến $p$ ta được một
xâu, ký hiệu $S_p$. Ngược lại, với một xâu $S$, ta đi theo các chữ cái của
$S$ trong trie và đến một nút cuối, ký hiệu $p_S$. Như vậy trong trie đầy đủ
có tương ứng một-một giữa nút và xâu.

\subsection{Xây dựng trie}

Trong thực tế ta không thể xây một trie vô hạn. Ta bắt đầu từ cây rỗng và lần
lượt chèn những từ cần dùng. Chẳng hạn, sau khi chèn \texttt{aa} rồi
\texttt{aza}, phần hữu hạn của trie chứa các đường:
\[
\begin{array}{c}
\varepsilon\\
\downarrow\texttt{a}\\
\texttt{a}\\
\swarrow\texttt{a}\quad\searrow\texttt{z}\\
\texttt{aa}\qquad\texttt{az}\\
\hfill\downarrow\texttt{a}\\
\hfill\texttt{aza}
\end{array}
\]

Quy trình chèn một xâu $S[1..n]$ là:

\begin{verbatim}
p <- Root
for i <- 1 to n do
  if p.Child[S[i]] = NIL then
    create and initialize p.Child[S[i]]
  p <- p.Child[S[i]]
\end{verbatim}

Ở đây \texttt{p.Child[ch]}, với \texttt{ch} thuộc \texttt{a..z}, là con của
$p$ theo cạnh mang nhãn \texttt{ch}.

\subsection{Định nghĩa con trỏ tiền tố}

Cây tiền tố từ khác trie thường ở chỗ mỗi nút không phải gốc có thêm một con
trỏ tiền tố. Với nút $p$, tìm số nguyên nhỏ nhất
\[
  k\in[2,|S_p|]
\]
sao cho $S_p[k..|S_p|]$ tồn tại trong cây. Khi đó con trỏ tiền tố của $p$ trỏ
đến nút
\[
  p_{S_p[k..|S_p|]}.
\]

Ví dụ, xét nút ứng với \texttt{abab}. Ta thử hậu tố \texttt{bab}; nếu
\texttt{bab} không xuất hiện trong trie, ta thử tiếp \texttt{ab}. Nếu
\texttt{ab} có trong trie, con trỏ tiền tố của \texttt{abab} trỏ đến nút
\texttt{ab}. Các nút rỗng không được biểu diễn trong các sơ đồ này; con trỏ
tiền tố nên hiểu như một cạnh phụ không thuộc cây.

\subsection{Sinh con trỏ tiền tố}

Nếu tính con trỏ tiền tố bằng cách dò lại từ định nghĩa, độ phức tạp sẽ rất
lớn. Bắt chước hàm tiền tố trong KMP, ta có thể dùng con trỏ tiền tố của nút
cha để suy ra con trỏ tiền tố của nút hiện tại. Gọi nút được con trỏ tiền tố
của $p$ trỏ tới là nút tiền tố của $p$.

Giả sử đang xử lý nút $p$ và mọi nút có độ sâu nhỏ hơn đã có con trỏ tiền tố.
Ta tính con trỏ của $p$ như sau:

\begin{verbatim}
q <- p.Father
ch <- label of the edge q -> p
q <- q.Prefix
while q != Root and q.Child[ch] = NIL do
  q <- q.Prefix
if q.Child[ch] = NIL then
  p.Prefix <- Root
else
  p.Prefix <- q.Child[ch]
\end{verbatim}

Nói bằng lời: từ nút tiền tố của cha, thử đi xuống bằng đúng chữ cái cuối cùng
dẫn vào $p$; nếu không có cạnh đó, tiếp tục đi theo các con trỏ tiền tố cho
đến khi tìm được hoặc về gốc.

\section{Thuật toán chính thứ nhất: thuật toán tuyến tính}

\subsection{Gợi ý từ KMP}

Tinh túy của KMP là giảm tính toán trùng lặp: dựa vào sự tự khớp sau khi dịch
của mẫu để xác định cần dịch mẫu bao nhiêu. Với bài toán nhiều mẫu, ta cũng
muốn chỉ quét văn bản một lần. Khi gặp ký tự không khớp, ta cần tìm hậu tố
dài nhất có thể của tiền tố hiện tại. Hậu tố ấy không nhất thiết phải là tiền
tố của cùng một mẫu; nó có thể là tiền tố của bất kỳ mẫu nào.

Ví dụ, văn bản có đoạn đầu \texttt{zababaa...}; mẫu \texttt{ababaca} thất bại
sau khi đã khớp \texttt{ababa}. Ta không nhất thiết dịch mẫu này sang phải
2 ô. Nếu có mẫu khác như \texttt{babaab}, ta có thể dịch ít hơn và tiếp tục từ
tiền tố của mẫu mới. Dịch càng ít càng tốt mới là điều cần thiết.

Các mẫu liên hệ với nhau khá chặt chẽ. Để hạ độ phức tạp cài đặt, cấu trúc
cây là một lựa chọn tự nhiên. Kết hợp với tư tưởng hàm tiền tố của KMP, ta thu
được cây tiền tố từ đã nêu ở Chương 3.

\subsection{Dùng cây tiền tố từ và nhãn phụ}

Ta xây cây tiền tố từ \texttt{TreeA} bằng cách chèn toàn bộ mẫu, rồi tính các
con trỏ tiền tố. Khi nào thì ta thật sự tìm thấy một mẫu $P_a$? Cách đơn giản
nhất là đặt nhãn phụ \texttt{Okay} tại nút $p_{P_a}$, với giá trị $a$, để ghi
nhận rằng đây là điểm kết thúc của mẫu $a$; các nút khác có
\texttt{Okay = 0}.

Chỉ đánh dấu như vậy vẫn chưa đủ. Nếu có hai mẫu \texttt{abcd} và \texttt{bc},
khi đang ở nút \texttt{abcd}, mẫu \texttt{bc} cũng vừa kết thúc, vì
\texttt{bc} là một hậu tố có mặt trong trie. Nút \texttt{bc} chính là nút tiền
tố của \texttt{abcd}. Vì vậy khi sinh con trỏ tiền tố, ta truyền nhãn
\texttt{Okay}: nếu nút tiền tố của nút hiện tại có nhãn \texttt{Okay}, nút
hiện tại cũng phải ghi nhận nhãn ấy. Trong cài đặt để tìm mọi mẫu, một nút nên
lưu danh sách các mẫu kết thúc tại nó thay vì chỉ một giá trị.

\subsection{Độ phức tạp của cây tiền tố từ}

Trong quá trình sinh con trỏ tiền tố, ta dùng đúng tư tưởng như khi sinh mảng
tiền tố KMP. Có thể chứng minh độ phức tạp bằng phân tích thế năng.

Với một mẫu $P[1..n]$, đặt $p_i$ là nút tương ứng với $P[1..i]$ trong trie,
và $p_0$ là gốc. Ta chứng minh tổng chi phí tính con trỏ tiền tố cho
$p_0,p_1,\ldots,p_n$ là $O(n)$.

Gọi $c_i$ là chi phí tính con trỏ tiền tố của $p_i$, với $c_0=1$. Định nghĩa
thế năng $\Phi(i)$ của $p_i$ như sau:
\[
  \Phi(0)=0,\qquad
  \Phi(i)=\text{độ sâu của nút tiền tố của }p_i\quad(1\le i\le n).
\]
Rõ ràng $p_i$ là cha của $p_{i+1}$ khi $0\le i<n$. Trong giả mã sinh con trỏ,
mỗi lần vòng lặp đi theo \texttt{q.Prefix}, độ sâu của $q$ giảm ít nhất $1$.
Do đó số lần lặp khi tính $p_{i+1}$ nhiều nhất là
\[
  \Phi(i)+1-\Phi(i+1).
\]
Suy ra
\[
  c_{i+1}\le \Phi(i)+2-\Phi(i+1).
\]
Cộng theo mọi $0\le i<n$:
\[
  \sum_{i=1}^{n} c_i
  \le 2n+\Phi(0)-\Phi(n)
  \le 2n.
\]
Vì vậy việc tính con trỏ tiền tố trên $n$ nút này mất $O(n)$ thời gian. Toàn
bộ cây tiền tố từ được xây trong
\[
  O\left(\sum L\right).
\]

\subsection{Quy trình chính}

Bắt chước KMP, thuật toán khớp nhiều xâu đầu tiên là:

\begin{verbatim}
Algorithm MultiPatternMatching1
  preprocess: build the prefix trie
  q <- Root
  for i <- 1 to n do
    while q != Root and q.Child[T[i]] = NIL do
      q <- q.Prefix
    if q.Child[T[i]] != NIL then
      q <- q.Child[T[i]]
    if q.Okay then
      output every pattern ending at i represented by q.Okay
End
\end{verbatim}

Giống KMP, số lần thực hiện \texttt{q <- q.Prefix} không vượt quá tổng số lần
đi xuống theo cạnh cây. Vì vậy phần quét văn bản chạy trong $O(n)$.

Nếu yêu cầu đúng vị trí bắt đầu nhỏ nhất của một mẫu, cần chú ý rằng mẫu kết
thúc sớm hơn chưa chắc bắt đầu sớm hơn. Ví dụ một mẫu dài hơn có thể bắt đầu
trước nhưng kết thúc sau một mẫu ngắn. Ta chỉ cần thêm phép so sánh vị trí bắt
đầu $i-L_a+1$ của mọi mẫu vừa được báo, rồi lấy vị trí nhỏ nhất.

\subsection{Độ phức tạp thời gian và bộ nhớ}

Chương trình chính gồm hai phần: xây cây tiền tố từ và quét văn bản. Tổng thời
gian là
\[
  O\left(\sum L+n\right).
\]

Nếu mỗi nút lưu đủ 26 con trỏ, mỗi nút cần $26$ ô con trỏ. Số nút là
$O(\sum L)$, do đó bộ nhớ là
\[
  O\left(26\sum L\right).
\]
Quét văn bản không cần thêm bộ nhớ đáng kể.

\subsection{Một số mở rộng}

Nếu văn bản không dài, hoặc bảng chữ cái lớn, bộ nhớ có thể trở thành nút thắt
cổ chai. Có hai hướng xử lý.

Thứ nhất, chuyển mọi ký tự thành nhị phân rồi xây cây nhị phân. Ví dụ chữ
\texttt{a} có thể được biểu diễn bằng một dãy bit cố định như
\texttt{01100001}. Khi đó cây 26 phân trở thành cây nhị phân. Mức cải thiện
không thật lý tưởng vì độ dài biểu diễn tăng lên.

Thứ hai, phân bổ động danh sách con của từng nút và lưu có thứ tự. Khi cần đi
theo một ký tự, ta tìm nhị phân trong danh sách con. Thời gian có tăng một
nhân tử nhỏ, nhưng bộ nhớ giảm còn
\[
  O\left(\sum L\right).
\]

Nếu yêu cầu mọi vị trí xuất hiện của mọi mẫu, thuật toán vẫn làm được. Điểm
khác là tại \texttt{Okay} có thể có nhiều mẫu cùng kết thúc, nên ta lưu một
danh sách thay vì một nhãn đơn.

Khi mẫu ngắn nhất rất dài và văn bản gần như ngẫu nhiên, còn có một thuật toán
có hiệu năng trung bình tốt hơn. Để trình bày nó, trước hết cần giới thiệu cây
hậu tố.

\section{Cây hậu tố và thuật toán McCreight}

\subsection{Cấu trúc dữ liệu}

Cây hậu tố cơ bản là trie tạo bởi tất cả hậu tố của một từ. Với từ
\texttt{ababc}, ta có năm hậu tố:
\[
  \texttt{ababc},\quad \texttt{babc},\quad
  \texttt{abc},\quad \texttt{bc},\quad \texttt{c}.
\]
Nếu xây trie trực tiếp, số nút có thể đạt cấp $O(N^2)$ với $N$ là độ dài từ.
Để tránh điều này, ta nén các đường chỉ có một nhánh: mỗi cạnh có thể mang cả
một xâu, và các cạnh đi ra từ cùng một nút có ký tự đầu khác nhau. Khi chèn
một hậu tố, nhiều nhất chỉ thêm một cạnh mới, nên số nút là $O(N)$.

Một vấn đề nhỏ là có thể có một hậu tố kết thúc ở giữa một cạnh đã nén. Để
tránh trường hợp ấy, ta thêm vào cuối xâu một ký tự chưa xuất hiện trong xâu,
ví dụ \texttt{\$}. Khi đó không hậu tố nào là tiền tố của hậu tố khác.

\subsection{Một số định nghĩa}

Cho từ $S[1..n]$. Ký hiệu
\[
  \operatorname{suf}_i=S[i..n]
\]
là hậu tố dài thứ $i$ của $S$.

Mỗi nút trong cây hậu tố có các trường chính:
\begin{itemize}
  \item \texttt{Father}: con trỏ đến cha;
  \item \texttt{Suffix}: liên kết hậu tố, sẽ định nghĩa bên dưới;
  \item \texttt{Child[c]}: con có cạnh bắt đầu bằng ký tự \texttt{c};
  \item \texttt{String[c]}: xâu ghi trên cạnh đến \texttt{Child[c]}.
\end{itemize}
Trong cài đặt, \texttt{String[c]} thường được biểu diễn bằng hai chỉ số đầu và
cuối trỏ vào xâu gốc $S$.

\subsection{Xây cây hậu tố: khung ban đầu}

Thuật toán McCreight xây cây hậu tố $T$ trong $n$ bước. Ban đầu $T$ rỗng; ở
bước $i$, chèn $\operatorname{suf}_i$ vào $T$.

Định nghĩa $\operatorname{head}_i$ là tiền tố chung dài nhất, lấy lớn nhất
trên mọi $\operatorname{suf}_j$ với $1\le j<i$, giữa $\operatorname{suf}_i$
và $\operatorname{suf}_j$. Viết
\[
  \operatorname{suf}_i=\operatorname{head}_i+\operatorname{tail}_i.
\]
Khung xây dựng là:

\begin{verbatim}
for i <- 1 to n do
  find the node corresponding to head[i]
  create p.Child[tail[i][1]] and initialize it
  p.String[tail[i][1]] <- tail[i]
\end{verbatim}

Ví dụ với $S=\texttt{abab\$}$, ta lần lượt chèn
\texttt{abab\$}, \texttt{bab\$}, \texttt{ab\$}, \texttt{b\$}, \texttt{\$}.
Việc thêm \texttt{tail[i]} chỉ tốn $O(1)$ nếu đã biết nút của
\texttt{head[i]}; nút này mới là phần khó. Nếu mỗi lần đều dò từ gốc, tổng
thời gian sẽ rất cao.

\subsection{Liên kết hậu tố}

Tính chất 1. Nếu $\operatorname{head}_{i-1}$ viết được thành $x\delta$, trong
đó $x$ là một ký tự và $\delta$ có thể rỗng, thì $\delta$ chắc chắn là tiền tố
của $\operatorname{head}_i$.

Thật vậy, cây rẽ nhánh tại nút ứng với $\operatorname{head}_{i-1}$, nên tồn
tại $j<i-1$ sao cho $x\delta$ là tiền tố của $\operatorname{suf}_j$. Khi đó
$\delta$ là tiền tố của $\operatorname{suf}_{j+1}$. Đồng thời $\delta$ cũng là
tiền tố của $\operatorname{suf}_i$, nên nó là tiền tố của
$\operatorname{head}_i$.

Liên kết hậu tố được định nghĩa như sau. Với mỗi nút trong không phải lá $p$,
xâu ứng với nó có thể viết là $x\delta$, với $x$ là một ký tự và $\delta$ có
thể rỗng. Con trỏ \texttt{p.Suffix} trỏ đến vị trí của $\delta$ trong cây; nếu
$\delta$ rỗng thì trỏ về gốc. Liên kết hậu tố của gốc trỏ đến chính gốc.

Trong quá trình xây, ta duy trì:

\begin{itemize}
  \item Tính chất 2: ngay sau khi một nút trong không phải lá được tạo, liên
  kết hậu tố của nó được tính ngay.
  \item Tính chất 3: khi cần tính liên kết hậu tố của một nút trong, xâu đích
  của liên kết hậu tố đã xuất hiện trong cây.
\end{itemize}

Nhờ vậy, tương tự cây tiền tố từ, để xây liên kết hậu tố của nút $p$ ta bắt
đầu từ liên kết hậu tố của cha $p$, rồi đi xuống:

\begin{verbatim}
Procedure CountSuffixLink(p)
  q <- p.Father.Suffix
  s <- string on the edge p.Father -> p
  if q = Root then
    s <- s[2..|s|]
  while s != "" do
    s1 <- q.String[s[1]]
    L <- min(|s|, |s1|)
    if L = |s| and L != |s1| then
      create an internal node r and initialize it
      r.String[s1[L+1]] <- s1[L+1..|s1|]
      r.Child[s1[L+1]] <- q.Child[s[1]]
      q.String[s[1]] <- s
      q.Child[s[1]] <- r
      p.Suffix <- r
      CountSuffixLink(r)
      return
    q <- q.Child[s[1]]
    s <- s[L+1..|s|]
End
\end{verbatim}

Vì tính chất 3 đảm bảo xâu cần tìm đã có trong cây, ở dòng tính $L$ chỉ cần so
sánh độ dài phần cạnh liên quan để biết đi tiếp hay tách cạnh.

\subsection{Xây cây hậu tố}

Khi đã có liên kết hậu tố, mỗi lần tìm $\operatorname{head}_i$ chỉ cần bắt đầu
từ liên kết hậu tố của nút ứng với $\operatorname{head}_{i-1}$ rồi đi xuống.
Khung thuật toán trở thành:

\begin{verbatim}
for i <- 1 to n do
  1. find the node p corresponding to head[i-1]
  2. from p.Suffix, descend to the node q for head[i]
  3. create q.Child[tail[i][1]] and initialize it
  4. q.String[tail[i][1]] <- tail[i]
  5. if p.Suffix != q then CountSuffixLink(q)
\end{verbatim}

Phân tích thời gian:
\begin{itemize}
  \item Các bước 1, 3 và 4 tốn $O(1)$ mỗi vòng, tổng cộng $O(N)$.
  \item Bước 5 tính liên kết hậu tố cho mọi nút trong. Vì trong
  \texttt{CountSuffixLink} ta dùng độ dài cạnh để đi qua hoặc tách cạnh, tổng
  chi phí tỉ lệ với số nút, tức $O(N)$.
  \item Bước 2 tốn $O(|\operatorname{head}_i|-|\operatorname{head}_{i-1}|+1)$
  theo phân tích khấu hao, tổng cộng cũng là $O(N)$.
\end{itemize}
Vì vậy thuật toán McCreight xây cây hậu tố trong $O(N)$.

\section{Thuật toán chính thứ hai: thuật toán có hiệu năng trung bình tốt hơn}

\subsection{Dùng và mở rộng cây tiền tố từ}

Tương tự thuật toán chính thứ nhất, ta dùng cây tiền tố từ \texttt{TreeA} xây
từ toàn bộ mẫu. Để phù hợp với thuật toán thứ hai, mỗi nút thêm tham số
\texttt{Shift}. Tham số này ghi độ dài đường đi ngắn nhất từ nút hiện tại đến
một nút \texttt{Okay} bất kỳ, không kể chính nó. Đường đi có thể dùng cạnh cây
hoặc con trỏ tiền tố; cạnh cây có trọng số $1$, con trỏ tiền tố có trọng số
$0$.

Ý nghĩa của \texttt{Shift}: nếu con trỏ hiện đang ở nút này, ta cần đọc thêm
ít nhất bao nhiêu ký tự nữa mới có thể tạo ra một lần khớp mẫu.

Với các mẫu \texttt{abab}, \texttt{ba}, \texttt{bb}, các nút kết thúc mẫu được
đánh dấu \texttt{Okay}. Giá trị \texttt{Shift} được xây như sau. Khi chèn ký
tự thứ $j$ của mẫu $P_i$, đặt
\[
  \operatorname{new}=
  \begin{cases}
    L_i-j, & 1\le j<L_i,\\
    L_i, & j=L_i.
  \end{cases}
\]
Nếu tạo nút mới $p$, gán \texttt{p.Shift <- new}. Nếu đi đến một nút đã tồn
tại $p$, gán
\[
  \texttt{p.Shift}\leftarrow
  \min(\texttt{p.Shift},\operatorname{new}).
\]

Như vậy mới xét các cạnh cây, chưa xét con trỏ tiền tố. Vì con trỏ tiền tố
luôn đi từ tầng sâu lên tầng nông hơn, ta duyệt các nút theo độ sâu tăng dần
và cập nhật
\[
  \texttt{p.Shift}\leftarrow
  \min(\texttt{p.Shift},\texttt{p.Prefix.Shift}).
\]
Tổng số nút là $O(\sum L)$, nên xây và duy trì \texttt{TreeA} mất
$O(\sum L)$ thời gian.

\subsection{Dùng và mở rộng cây hậu tố}

Thuật toán thứ hai còn dùng một cây hậu tố \texttt{TreeB}, được xây từ tất cả
hậu tố của các mẫu sau khi đảo ngược.

Nếu các mẫu là \texttt{abab}, \texttt{ba}, \texttt{bb}, các xâu đảo lần lượt
là \texttt{baba}, \texttt{ab}, \texttt{bb}. Cây hậu tố của chúng cho phép ta
kiểm tra từ phải sang trái một đoạn văn bản và biết trong $O(N)$ thời gian
liệu đoạn độ dài $N$ ấy có phải là xâu con của một mẫu nào không.

Khác phần trình bày cây hậu tố ở trên, trong mô tả \texttt{TreeB} ta không cần
hiện các cạnh chứa ký tự kết thúc \texttt{\$}; ký tự ấy chỉ phục vụ quá trình
xây dựng.

\subsection{Hai hàm trên \texttt{TreeA} và \texttt{TreeB}}

\begin{verbatim}
Function ScanA(L, R, point)
  scan T[L..R] in TreeA, starting from node point
  set point to the last reached node
  while point.Shift < shortest_pattern_length div 2 do
    read one more character from the text
    continue moving point in TreeA
  output every Okay encountered during the process
  return the last text position scanned and the final point
End
\end{verbatim}

Quét trong \texttt{TreeA} chính là bước đi của thuật toán ở Chương 4: đi xuống
theo cạnh cây khi có thể, nếu không thì theo con trỏ tiền tố để lùi. Theo phân
tích khấu hao, thời gian tỉ lệ với số ký tự đã quét. Để trình bày gọn, từ đây
ta xét bài toán tìm mọi vị trí xuất hiện của mọi mẫu.

\begin{verbatim}
Function ScanB(L, R)
  scan T[L..R] from right to left in TreeB
  stop when no path matches, or when reaching L
  return the last text position scanned
End
\end{verbatim}

Nếu hàm này trả về $L$, đoạn $T[L..R]$ là xâu con của một mẫu nào đó. Ta gọi
đoạn như vậy là \term{hợp lệ}; ngược lại là \term{không hợp lệ}.

\subsection{Quy trình chính}

\begin{verbatim}
Algorithm MultiPatternMatching2
  preprocess: build TreeA and TreeB
  L <- 0
  R <- shortest_pattern_length
  point <- TreeA.Root
  while R <= n do
    pos <- ScanB(L+1, R)
    if pos > L + 1 then
      point <- TreeA.Root
    (L, point) <- ScanA(pos, R, point)
    R <- L + point.Shift
End
\end{verbatim}

Tư tưởng cơ bản là như sau. Điểm bắt đầu đang xét là $L+1$. Nếu tại đó có một
mẫu xuất hiện, đuôi của mẫu phải nằm sau $R$. Ta quét ngược từ $R$ bằng
\texttt{ScanB} và giả sử nó dừng ở \texttt{pos}.

Nếu \texttt{pos = L+1}, đoạn $T[L+1..R]$ là hợp lệ, nên có thể tiếp tục quét
tiến bằng \texttt{ScanA} từ $L+1$.

Nếu \texttt{pos > L+1}, thì mọi vị trí $L+1,L+2,\ldots,\texttt{pos}-1$ đều
không thể là đầu của một lần khớp mẫu. Thật vậy, nếu tại một vị trí
$x$ trong khoảng đó có mẫu xuất hiện, đoạn $T[x..R]$ phải là hợp lệ, mâu
thuẫn với việc \texttt{ScanB} dừng tại \texttt{pos}. Đoạn $T[\texttt{pos}..R]$
là hợp lệ, nên ta đặt lại \texttt{point} về gốc của \texttt{TreeA} rồi quét
tiến từ \texttt{pos}.

Điều kiện dừng của \texttt{ScanA},
\[
  \texttt{point.Shift}\ge \left\lfloor\frac{\theta}{2}\right\rfloor,
\]
với $\theta$ là độ dài mẫu ngắn nhất, nhằm bảo đảm ở giai đoạn kế tiếp
$R-L$ có cỡ ít nhất $\theta/2$.

\subsection{Một ví dụ}

Lấy hai mẫu \texttt{abcd} và \texttt{bcde}. Cây tiền tố \texttt{TreeA} gồm hai
đường:
\[
\begin{array}{cccccc}
1&\xrightarrow{\texttt{a}}&2&\xrightarrow{\texttt{b}}&3&
\xrightarrow{\texttt{c}}4\xrightarrow{\texttt{d}}5,\\[0.4em]
1&\xrightarrow{\texttt{b}}&6&\xrightarrow{\texttt{c}}&7&
\xrightarrow{\texttt{d}}8\xrightarrow{\texttt{e}}9.
\end{array}
\]
Các nút $5$ và $9$ là nút \texttt{Okay}. Bảng \texttt{Shift} là:
\[
\begin{array}{c|ccccccccc}
\text{nút} & 1&2&3&4&5&6&7&8&9\\
\hline
\texttt{Shift} & 4&3&2&1&1&3&2&1&4
\end{array}
\]

Giả sử văn bản $T=\texttt{abcabcde}$.

Giai đoạn 1: $L=0$, $R=4$, \texttt{point = 1}. Quét ngược bằng \texttt{ScanB}
từ $R$ về $L$. Chỉ đoạn \texttt{a} là hợp lệ, còn \texttt{ca} không hợp lệ,
nên \texttt{pos = 4}. Vì \texttt{pos > L+1}, các vị trí $1..3$ không thể bắt
đầu một mẫu; nếu có, mẫu ấy phải chứa đoạn \texttt{ca}, nhưng \texttt{ca}
không là xâu con của mẫu nào. Sau đó \texttt{ScanA} quét đoạn
$T[4..4]=\texttt{a}$, đưa \texttt{point} từ $1$ đến $2$. Vì
\texttt{Shift[2] = 3} không nhỏ hơn $\lfloor 4/2\rfloor$, giai đoạn kết thúc
và chưa tìm thấy mẫu.

Giai đoạn 2: $L=4$, $R=L+\texttt{Shift[point]}=7$, \texttt{point = 2}. Quét
ngược thấy \texttt{bcd} hợp lệ và dừng tại \texttt{pos = 5 = L+1}, nên không
cần đặt lại \texttt{point}. Quét tiến $T[5..7]$ đưa \texttt{point} đến nút
$5$ và phát hiện mẫu \texttt{abcd}. Vì $\texttt{Shift[5]} = 1 <
\lfloor4/2\rfloor$, \texttt{ScanA} đọc thêm ký tự \texttt{e}, đi qua các con
trỏ và cạnh tương ứng đến nút $9$, rồi phát hiện tiếp mẫu \texttt{bcde}. Khi
hết văn bản, thuật toán dừng.

\subsection{Phân tích độ phức tạp}

Xây cây tiền tố từ và cây hậu tố tốn
\[
  O\left(26\sum L\right),
\]
trong đó bộ nhớ vẫn có thể là nút thắt cổ chai như đã nêu ở Chương 4. Ta tập
trung phân tích quá trình quét văn bản.

Gọi $\theta$ là độ dài mẫu ngắn nhất. Trường hợp xấu nhất, \texttt{ScanA} có
thể phải quét toàn bộ văn bản vì \texttt{Shift} luôn nhỏ hơn $\theta/2$. Khi
đó thuật toán vẫn tốn $O(n)$ thời gian.

Khi mọi mẫu đều có độ dài $\theta$, $\theta$ đủ lớn, văn bản ngẫu nhiên và
bảng chữ cái đủ lớn, xác suất để \texttt{ScanB} quét hết $T[L+1..R]$ không
lớn. Khi ấy độ phức tạp trung bình nhỏ hơn nhiều so với $O(n)$. Ta tính như
sau.

Giả sử mọi mẫu có độ dài $\theta$ và
\[
  M=\sum L=\theta^k,
\]
trong đó $k$ là hằng số nhỏ vì $\theta$ lớn. Gọi $\alpha$ là kích thước bảng
chữ cái, chẳng hạn $26$. Với $\theta$ và $\alpha$ đủ lớn, giả sử
\[
  \log_\alpha M^3 \le \frac{\theta}{2},
  \qquad\text{tức}\qquad
  3k\log_\alpha\theta\le \frac{\theta}{2}.
\]
Một lần \texttt{ScanB} cộng với một lần \texttt{ScanA} được gọi là một giai
đoạn.

\paragraph{Mệnh đề 1.}
Trong mỗi giai đoạn, $R-L$ có bậc $O(\theta)$.

Lần đầu vào vòng lặp, $L=0$ và $R=\theta$. Các lần sau, vì \texttt{ScanA} chỉ
dừng khi $\texttt{Shift}\ge \theta/2$, ta có $R\ge L+\theta/2$. Do đó
$R-L$ luôn có cỡ $O(\theta)$ trong phân tích này.

\paragraph{Mệnh đề 2.}
Số giai đoạn nhiều nhất là $O(n/\theta)$. Điều này suy ra trực tiếp từ mệnh
đề 1.

\paragraph{Mệnh đề 3.}
Tập mẫu $P$ có không quá $\theta^{2k}$ xâu con, tức không quá chừng ấy xâu hợp
lệ.

Mỗi mẫu có không quá $\theta^2$ xâu con. Số mẫu là $\theta^{k-1}$, nên tổng số
xâu con không vượt quá một cận cùng bậc; trong lập luận của tác giả, dùng cận
rộng $\theta^{2k}$ là đủ.

\paragraph{Mệnh đề 4.}
Tồn tại hằng số $C$ sao cho một xâu gồm $C\log_\alpha\theta$ ký tự ngẫu nhiên
có xác suất không quá $1/M$ để là xâu con của một mẫu trong $P$.

Theo mệnh đề 3, số xâu hợp lệ không quá $\theta^{2k}$. Số xâu độ dài
$C\log_\alpha\theta$ là
\[
  \alpha^{C\log_\alpha\theta}=\theta^C.
\]
Vì các khả năng được đọc vào có xác suất như nhau, xác suất đọc được xâu hợp
lệ không vượt quá
\[
  \frac{\theta^{2k}}{\theta^C}.
\]
Muốn giá trị này không quá $1/M=1/\theta^k$, chỉ cần $C\ge 3k$. Chọn $C=3k$
là một hằng số; giả thiết ở trên cũng cho
$C\log_\alpha\theta\le\theta/2$.

\paragraph{Mệnh đề 5.}
Trong \texttt{ScanA}, khi xử lý đến nút $p$ ở độ sâu $x$ trên \texttt{TreeA}
và tương ứng với ký tự văn bản $T[y]$, đoạn $T[y-x+1..y]$ là hợp lệ.

Thật vậy, đoạn này chính là xâu tạo bởi đường đi từ gốc \texttt{TreeA} đến
$p$. Nó là tiền tố của một mẫu, do đó là xâu con của một mẫu.

\paragraph{Mệnh đề 6.}
Trong \texttt{ScanA}, số ký tự kỳ vọng cần đọc thêm ngoài $[L..R]$ cho đến khi
gặp nút có \texttt{Shift} ít nhất $\theta/2$ là
\[
  O(\log_\alpha\theta).
\]

Nếu \texttt{ScanA} dừng trong $C\log_\alpha\theta$ ký tự đầu, kết quả hiển
nhiên. Nếu nó dừng trong khoảng từ $C\log_\alpha\theta$ đến
$2C\log_\alpha\theta$, thì sau $C\log_\alpha\theta$ ký tự đầu, con trỏ trong
\texttt{TreeA} phải có độ sâu ít nhất $\theta/2$; nếu độ sâu nhỏ hơn, vì
$C\log_\alpha\theta\le\theta/2$, giá trị \texttt{Shift} đã đủ lớn để dừng.
Theo mệnh đề 5, xâu dài $C\log_\alpha\theta$ vừa tạo là hợp lệ, xác suất của
sự kiện này không quá $1/M$ theo mệnh đề 4.

Tương tự, xác suất để dừng trong khoảng
$xC\log_\alpha\theta$ đến $(x+1)C\log_\alpha\theta$ giảm theo lũy thừa của
$1/M$. Vì vậy kỳ vọng được chặn bởi một tổng hình học:
\[
  \sum_{x\ge 0}
  (x+1)C\log_\alpha\theta\cdot \frac{1}{M^x}
  =O(\log_\alpha\theta).
\]

\paragraph{Mệnh đề 7.}
Số ký tự kỳ vọng được đọc trong mỗi giai đoạn là
\[
  O(\log_\alpha\theta).
\]

Chia làm hai trường hợp.

Trường hợp 1: \texttt{ScanB} đọc không quá $C\log_\alpha\theta$ ký tự. Theo
mệnh đề 4, xác suất của trường hợp này ít nhất $1-1/M$. Khi đó \texttt{ScanA}
bắt đầu từ gốc của \texttt{TreeA} và đọc nhiều nhất
$C\log_\alpha\theta$ ký tự nữa; vì
$C\log_\alpha\theta\le\theta/2$, tham số \texttt{Shift} chắc chắn đạt
$\theta/2$. Tổng số ký tự không quá $2C\log_\alpha\theta$.

Trường hợp 2: \texttt{ScanB} vẫn chưa dừng sau
$C\log_\alpha\theta$ ký tự. Xác suất không vượt quá $1/M$. Khi đó kỳ vọng số
ký tự đọc vào không vượt quá $\theta$ cho \texttt{ScanB}, $\theta$ cho phần
\texttt{ScanA} bên trong đoạn, cộng thêm $C\log_\alpha\theta$ ký tự đọc thêm
do điều kiện \texttt{Shift}.

Do đó số ký tự kỳ vọng trong một giai đoạn là
\[
  \left(1-\frac1M\right)2C\log_\alpha\theta
  +\frac1M\left(2\theta+C\log_\alpha\theta\right)
  =O(\log_\alpha\theta).
\]

Kết hợp mệnh đề 2 và mệnh đề 7, độ phức tạp trung bình của thuật toán là
\[
  O\left(\frac{n\log_\alpha\theta}{\theta}\right).
\]
Hằng số ẩn khá lớn, nên ưu thế của thuật toán chỉ thể hiện rõ khi $\theta$ và
$\alpha$ đủ lớn.

\section{Gợi ý và kết luận}

\subsection{Phân tích thuật toán}

Trong thuật toán chính thứ hai, số $\theta/2$ mang hai ý nghĩa cùng lúc:
\begin{itemize}
  \item nó là ngưỡng để \texttt{ScanA} dừng;
  \item nó là cận dưới cho độ dài $R-L$ trong mỗi giai đoạn.
\end{itemize}

Nếu tăng ngưỡng này, \texttt{ScanA} khó dừng hơn. Nếu giảm nó, số giai đoạn sẽ
tăng. Cả hai đều làm độ phức tạp trung bình tăng. Vì vậy $\theta/2$ là điểm
trung gian vừa đủ để đạt độ phức tạp trung bình nhỏ nhất trong phân tích này.

\subsection{Gợi ý phương pháp}

Từ thuật toán này có thể rút ra hai gợi ý:
\begin{itemize}
  \item phối hợp nhiều thuật toán hoặc cấu trúc dữ liệu, như cây tiền tố từ và
  cây hậu tố trong bài này;
  \item chọn điểm cân bằng thích hợp giữa hai yếu tố đối nghịch.
\end{itemize}

Trong bài này ta dùng trung bình cộng để chọn điểm cân bằng. Thường gặp hơn là
trung bình nhân. Một ví dụ điển hình là bài \emph{Editor} ở ngày thi thứ nhất
NOI 2003: trong một mảng ký tự dài tối đa $2M$, liên tục thực hiện chèn, xóa,
định vị, đọc, v.v., và cần mô phỏng quá trình đó.

Nếu dùng mảng, chèn và xóa tốn $O(2M)$, còn định vị tốn $O(1)$. Nếu dùng danh
sách liên kết, chèn và xóa tốn $O(1)$, còn định vị tốn $O(2M)$. Cả hai cách
đều không đủ mạnh với giới hạn dữ liệu. Ta phối hợp chúng bằng danh sách liên
kết theo khối:
\[
\boxed{\text{mảng khối 1}}\to
\boxed{\text{mảng khối 2}}\to
\cdots\to
\boxed{\text{mảng khối }Y}.
\]
Mỗi khối là một mảng. Đặt kích thước khối là $X$ và số khối là $Y$; lý tưởng
thì $XY$ xấp xỉ độ dài văn bản. Khi chèn hoặc xóa, ta thao tác vụn trong một
vài khối với chi phí $O(X)$, rồi chỉnh con trỏ danh sách với chi phí $O(1)$.
Khi định vị, ta tìm khối với chi phí $O(Y)$ rồi tìm vị trí trong khối với chi
phí $O(1)$.

Chọn $X=Y=\sqrt{2M}$ tạo điểm cân bằng và hạ độ phức tạp tổng thể.

\subsection{Kết luận}

Việc cài đặt thuật toán chính thứ hai đã vượt xa độ khó thông thường của OI;
hơn nữa thời gian đọc mảng ký tự có thể đã lớn hơn thời gian xử lý chính. Khi
đã có thuật toán chính thứ nhất đơn giản và dễ dùng, thuật toán thứ hai còn ý
nghĩa thực tế không?

Có. Trước hết, các thuật toán chuẩn bị được dùng ở đây đều đã trở thành công
cụ chủ lưu trong thi tin học; nắm vững chúng là năng lực cần có. Tiếp theo,
độ phức tạp của thuật toán thứ hai phản ánh ở một mức độ khác rằng cán cân
mâu thuẫn giữa phần cứng và phần mềm lại nghiêng về phía phần mềm: phương
tiện lưu trữ và thuật toán lưu trữ hiệu quả hơn vẫn cần được tiếp tục phát
triển.

Điều đáng quý nhất là tư tưởng ``dùng tổng hợp, mở từ điểm yếu nhất'' trong
cách hình thành thuật toán. Độ bền của một chuỗi phụ thuộc vào mắt yếu nhất;
lượng nước trong thùng phụ thuộc vào tấm ván ngắn nhất. Chỉ cần liên tục đột
phá ở nút thắt, ta sẽ luôn tìm được đường đi qua bài toán.

Tin học là một ngành sâu rộng. Bản chất của nó không nằm ở kỹ thuật lập trình
tầng thấp, mà nằm ở đổi mới tư duy của con người; đó cũng là mục đích của các
kỳ thi học thuật.

\section*{Lời cảm ơn}

Tác giả cảm ơn toàn thể các thầy cô trong nhóm huấn luyện của Ủy ban Olympic
Tin học thanh thiếu niên tỉnh Giang Tô; cảm ơn các thầy cô của Trường Ngoại
ngữ Nam Kinh đã hỗ trợ; cảm ơn bạn Xin Tao ở Liêu Ninh đã giúp đỡ nhiều; cảm
ơn gia đình đã quan tâm; và cảm ơn tất cả bạn bè trong cộng đồng tin học.

\begin{thebibliography}{9}
\bibitem{clrs}
Thomas H. Cormen, Charles E. Leiserson, Ronald L. Rivest, Clifford Stein,
\emph{Introduction to Algorithms}.

\bibitem{taocp}
Donald E. Knuth,
\emph{The Art of Computer Programming}.

\bibitem{csehandbook}
Allen B. Tucker et al.,
\emph{The Computer Science and Engineering Handbook}.

\bibitem{ukkonen}
E. Ukkonen,
\emph{On-line Construction of Suffix Trees}, 1995.

\bibitem{mccreight}
Barbara König,
\emph{McCreight's Suffix Tree Construction Algorithm}.

\bibitem{practical}
Wu Wenhu and Wang Jiande,
\emph{Analysis and Programming of Practical Algorithms}.
\end{thebibliography}

\end{document}
